\documentclass[10pt,a4paper]{article}
\usepackage[margin=1.55cm]{geometry}
\usepackage{graphicx}
\usepackage{xcolor}
\usepackage{enumitem}
\usepackage{tabularx}
\usepackage{array}
\usepackage{titlesec}
\usepackage{microtype}
\usepackage[hidelinks]{hyperref}
\usepackage{fontspec}
\definecolor{ink}{HTML}{17202A}
\definecolor{muted}{HTML}{5E6B75}
\definecolor{accent}{HTML}{1D6B66}
\definecolor{rule}{HTML}{D9E0DD}
\pagestyle{empty}
\setlength{\parindent}{0pt}
\setlength{\parskip}{0pt}
\setlist[itemize]{leftmargin=1.2em,itemsep=1pt,topsep=2pt,parsep=0pt}
\titleformat{\section}{\large\bfseries\color{accent}}{}{0pt}{}[\vspace{2pt}\color{rule}\titlerule]
\titlespacing*{\section}{0pt}{10pt}{5pt}
\newcommand{\entry}[4]{%
  \begin{tabularx}{\textwidth}{@{}p{2.45cm}@{\hspace{0.28cm}}X@{}}
  \textcolor{muted}{#1} & \textbf{#2}\\[-1pt]
  & #3\\[-1pt]
  & {\small\textcolor{muted}{#4}}
  \end{tabularx}\vspace{4pt}}
\newcommand{\paper}[4]{%
  \begin{tabularx}{\textwidth}{@{}p{1cm}@{\hspace{0.15cm}}X@{}}
  \textcolor{accent}{#1} & \textbf{#2}\\[-1pt]
  & #3\\[-1pt]
  & {\small\textcolor{muted}{#4}}
  \end{tabularx}\vspace{5pt}}
\begin{document}
\begin{minipage}[t]{0.78\textwidth}
  {\fontsize{26}{29}\selectfont\bfseries\color{ink} Junze Zhu}\par
  \vspace{7pt}{\color{accent}\bfseries Ph.D. Candidate \quad | \quad Artificial Intelligence, Nanjing University}\par
  \vspace{6pt}\small
  \href{mailto:aminataconde675@gmail.com}{aminataconde675@gmail.com}\quad|\quad
  \href{https://huicir.github.io/}{huicir.github.io}\quad|\quad
  \href{https://scholar.google.com/citations?user=-XbdENoAAAAJ\&hl=en}{Google Scholar}\quad|\quad
  \href{https://github.com/HuiCir}{GitHub}
\end{minipage}\hfill
\begin{minipage}[t]{0.18\textwidth}
  \vspace{-2pt}\IfFileExists{profile.jpg}{\includegraphics[width=2.3cm,height=2.8cm,keepaspectratio]{profile.jpg}}{}
\end{minipage}

\section{Research interests}
Language agents, multi-agent systems, speculative and asynchronous execution, harness-aligned policy learning, knowledge graphs, and recommendation.

\section{Education}
\entry{2025--2029 (expected)}{Ph.D. Candidate in Artificial Intelligence}{Nanjing University}{Focus: Agent Systems and Swarm Intelligence}
\entry{2022--2025}{Master's in Data Science}{Wuhan University}{Focus: Knowledge-Enhanced Intelligent Systems}
\entry{2018--2025}{Bachelor's in Electronic Information}{Wuhan University}{Focus: Multimodal Fusion Perception}

\section{Selected publications}
\paper{2026}{Recognize Your Orchestrator: An Entropy Dynamics Perspective for LLM Multi-Agent Systems}{Junze Zhu, Weihao Chen, Xuanwang Zhang, Zhen Wu, Xinyu Dai. \textit{ICML 2026}, PMLR 306:166844--166869. \href{https://proceedings.mlr.press/v306/zhu26g.html}{Paper}.}{Mean-field entropy dynamics for understanding instability in centralized language-agent orchestration.}
\paper{2025}{Knowledge Graph Preference Contrastive Learning for Recommendation}{Junze Zhu, Zhongyi Hu, Fan Zhang. \textit{Learning on Graphs Conference}, PMLR 269:18:1--18:15. \href{https://proceedings.mlr.press/v269/zhu25a.html}{Paper} / \href{https://github.com/HuiCir/KPCL}{Code}.}{Preference-aware contrastive learning and rationale attention for knowledge-aware recommendation.}
\paper{2026}{WebProgramBench: Can Coding Agents Reconstruct Live Applications from Scratch via Black-Box Interaction?}{Xuanwang Zhang, Yuteng Han, Xinyu Liu, \textbf{Junze Zhu}, et al. \textit{ICLR 2027 submission}, under review; fourth author. \href{https://openreview.net/forum?id=RKsIvqqtT8}{OpenReview}.}{A benchmark for executable reconstruction of live applications through black-box interaction.}

\section{Manuscripts and research projects}
\begin{itemize}
  \item \textbf{Anchoring and Adjustment: A Cognitive Refinement Framework for Trustworthy Retrieval Augmentation}. Working manuscript, 2026.
  \item \textbf{Perseus: Non-Blocking Parallel Exploration with a Speculative Swarm}. Under review. An asynchronous agent swarm architecture that expands inference-time exploration while an authoritative Actor retains task control. \href{https://huicir.github.io/Perseus/}{Project site}.
  \item \textbf{Harness Policy Distillation: Aligning Agent Policies with Runtime Semantics}. Under review. A policy learning framework that distills verified actions and runtime control-document views into a harness-compatible agent policy.
\end{itemize}

\section{Selected systems work}
\begin{itemize}
  \item \textbf{Open Crowd} -- skill and CLI management for complex task orchestration with agent swarms.
  \item \textbf{Token-Never-Sleeps} -- long-horizon agent swarm task management and external-loop execution; published as an npm package.
\end{itemize}

\section{Links}
\small
\href{https://huicir.github.io/}{Academic homepage}\quad|\quad
\href{https://github.com/HuiCir}{GitHub}\quad|\quad
\href{https://scholar.google.com/citations?user=-XbdENoAAAAJ\&hl=en}{Google Scholar}\quad|\quad
\href{https://openreview.net/forum?id=RKsIvqqtT8}{WebProgramBench on OpenReview}

\vfill
\begin{center}\scriptsize\color{muted}Last updated October 2026\end{center}
\end{document}
